\documentclass[fandol]{ctexart}
\usepackage{amsmath}
\usepackage{zhlineskip}
\pagestyle{empty}
\begin{document}
一般书强调“克氏符不是张量”是对的，本书则应强调
\begin{equation}
  [u, v](f) = u(v(f)) - v(u(f)),
\end{equation}
“克氏符是坐标系依赖的张量”。
\begin{equation}
  \begin{split}
    [u, v](f)
    &= u(v(f)) - v(u(f)) \\
    &= u^b \nabla_b (v^a \nabla_a f) - v^b \nabla_b (u^a \nabla_a f). \\
  \end{split}
\end{equation}
读者常问：你们为什么非把克氏符说成张量不可？
\end{document}
